\documentclass[12pt]{article}
\usepackage[utf8]{inputenc}
\usepackage[T1]{fontenc}
\usepackage{hyperref}
\title{Chingoni Language Resources for K-12 Curriculum Development}
\author{Compiled by Claude Sonnet 5.5 (claude-sonnet-5-5) for Andr\'as Kornai}
\date{}
\begin{document}
\maketitle

\section{Introduction}
Chingoni (Ngoni, Xingoni; ISO 639-3 xnj; Glottocode tanz1241) is a Bantu language spoken mainly around Songea in southwestern Tanzania (Ruvuma, Iringa and Morogoro regions are listed by Omniglot). Wikipedia says it is not a Nguni language despite its name, while another source places it in the Nguni subgroup, so classification is disputed. Speaker estimates range from about 305,000 to 313,000. The Tanzanian variety differs from the Mozambican one, with only marginal mutual intelligibility. The search found only general reference pages, and no school materials.

\section{Primary Educational Resources}
No primers, textbooks or school programmes for Chingoni were found. Swahili is the language of schooling in Tanzania, and Joshua Project notes interaction with Swahili and Makhuwa.

\section{Linguistic and Reference Materials}
\begin{itemize}
\item Wikipedia, Ngoni language (Tanzania): \url{https://en.wikipedia.org/wiki/Ngoni_language_(Tanzania)}
\item ISO 639 overview for Ngoni: \url{https://wikipedia.com/wiki/ISO_639:ngo}
\item Omniglot, Ngoni (Chingoni / Xingoni): \url{https://omniglot.com/writing/ngoni.htm}
\item Omniglot, Ngoni numbers: \url{https://omniglot.com/language/numbers/ngoni.htm}
\item Grambank entry: \url{https://grambank.clld.org/languages/tanz1241}
\end{itemize}

\section{Religious and Community Resources}
Joshua Project people-group profile, which says most Chingoni are Muslim. Its figures are less reliable than the linguistic sources: \url{https://joshuaproject.net/people_groups/13367}. No community organisations or scripture translations were confirmed in the search.

\section{Digital and Audio Resources}
No audio, corpora or software for Chingoni were found in the search.

\section{Curriculum Development Context}
Chingoni is widely used in its community, which is a strong basis for a mother-tongue programme. The search found no existing orthography standard, textbooks or teacher training. Next steps are to contact Tanzanian educational authorities and the Ruvuma/Songea community, to document an orthography, and to separate the Tanzanian variety from Mozambican Ngoni.
\end{document}